\scanpage{303}
\mspage{2}
Un objet de $A_*$ sera généralement noté par un symbole de la forme $K_*$, où $K_*$ désigne la famille des
\[ K_*(\Delta^n)=K_{[n]} \]
avec les opérations $\tfrac12$ simpliciales entre elles.
\margin{par abus de notation $K_*=(K_{[n]})_{n\geq0}$ ou mieux $\Delta^n\longmapsto K_{[n]}$}
On fera attention qu'on écrit $K_{[n]}$ et non $K_n$, pour des raisons qui vont apparaître (impérieuses lorsque $A$ est additive\ldots).
De même un objet de $A^*$ sera noté
\[ K^*=(n\longmapsto K^{[n]})\overset{\text{abus.}}{=}(K^{[n]})_n \]

Nous aurons à travailler avec la situation, où on a deux catégories $A$ et $B$, et une équivalence (notée \deleted{$\unsure{\vee^*}$} $M\mapsto M^\vee$)
\[ \vee:A^o\overset\simeq\longrightarrow B\qquad\text{d'où }B^o\simeq A \]
[le plus souvent $A=B$ et $(M^\vee)^\vee\underset{\substack{\text{isom}\\\text{fonct.}}}{\simeq}M$, avec compatibilités habituelles d'une autodualité\ldots]
on notera alors \inserted{souvent} par la \unsure{m\textsuperscript{e}} lettre un objet de $A_*$ (ou $A^*$) et l'objet de $A^*$ (resp. $A_*$) qui lui correspond par application de $\vee$, mais en indiquant la variance par la position du \uline{signe} $*$
\[
\begin{array}{ll}
K^*=(K_*)^\vee&K^{[n]}=K_{[n]}^\vee\\
K_*=(K^*)^\vee&K_{[n]}=K^{[n]\vee}
\end{array}
\]

On s'intéressera surtout au cas où ($h$ étant un anneau fixé) \deleted{$\unsure{A=\mathcal L}$}
\deleted{$\unsure{A:\ A}$}
\[
\begin{aligned}
A&=\deleted{\unsure{\mathrm{Ab}_h}}\ h\text{-modules à g. proj. d.t.f.}\\
B&=h\text{-modules à droite proj. de t.f.}
\end{aligned}
\]
(si $h$ est commutatif, on a $A=B$ et on a une autodualité dans $A$).

On a le foncteur canonique pl. fid
\diagram{\begin{tikzcd}(*)\quad\boldsymbol\Delta\arrow[r,hook]&\widehat{\boldsymbol\Delta}=\mathrm{Ss}\end{tikzcd}}
l'image par ce foncteur de $\Delta^n$ est noté $\Delta^n_*$ \inserted{ou mieux $\Delta^{[n]}_*$ (puisqu'il est covariant en $n$)}. On a donc
\[
\Delta^n_*\ (\text{ou }\Delta^{[n]}_*)=(m\longmapsto\Delta^n_{[m]}=\operatorname{Hom}_{\mathrm{Ens}}(\Delta^m,\Delta^n))
\]
et pour $\unsure{X}_*\in\operatorname{Ob}\widehat{\boldsymbol\Delta}$,
\[
\operatorname{Hom}(\Delta^n_*,\unsure{X}_*)\simeq X_*(\Delta^n)\simeq X_{[n]}
\]
en particulier
\[
\operatorname{Hom}(\Delta^n_*,\Delta^m_*)\simeq\Delta^m_*(\Delta^n_*)\simeq\operatorname{Hom}_{\boldsymbol\Delta}(\Delta^n,\Delta^m)
\]
\margin{$\Vert\quad(\operatorname{Hom}_{\mathrm{ord}}(\Delta^n,\Delta^m))$}
Pour $n$ variable, $\Delta^n_*\in\widehat{\boldsymbol\Delta}$ dépend de façon \uline{covariante} de $n$, \deleted{\unsure{(i.e. on}} donc le foncteur $n\longmapsto\Delta^n_*$ (qui n'est autre que (*)) est un objet canonique de $\widehat{\boldsymbol\Delta}^*$, qu'on appellera \uline{objet cosimplicial universel} et noté $\Delta^*$ (ou mieux $\Delta^*_*${)}
\[
\Delta^*_*\in\widehat{\boldsymbol\Delta}^*=(\mathrm{Ens}_*)^*\qquad\text{donc }\forall n\geq0
\]
\[
\Delta^{[n]}_*\in\widehat{\boldsymbol\Delta}\qquad
\Delta^{[n]}_*=\Delta^n_*=m\longmapsto
\underset{\substack{\Vert\\\operatorname{Hom}_{\boldsymbol\Delta}(\Delta^n,\Delta^m)}}{\Delta^n_{[m]}}
\]
\deleted{\unsure{Elle}}
[Une parenthèse de notations : Considérons les catégories de la forme $(A_*)^*$, un objet de \unsure{là-dedans} est donc noté $K^*=(K^{[n]})_n$, où les $K^{[n]}$ sont des $\in A_*$, donc s'écrivent $K^{[n]}_*$,
\[ K^{[n]}_*=(K^{[n]}_{[m]})_{m\geq0} \]
Ainsi, on est amené à noter $K^*_* $ les objets de $(A_*)^*\simeq\underline{\operatorname{Hom}}(\boldsymbol\Delta\times\boldsymbol\Delta^o,A$
\deleted{\unsure{on}}

\margin{304}
